% !TEX root = chapter-standalone.tex

\section{Initial objects}
\label{sec:initional}

The second ``role'' that we consider is that of ``being an initial object''. It is ``dual'' to the role of being a terminal object in that it is analogous to the latter, but with the direction of morphisms reversed. 


\begin{ctdefinition}\label{def:initial-object}
    Let $\CatC$ be a category.
    An \emph{initial object} in $\CatC$ is:

    \constit
    \begin{enumerate}
        \item an object $\initobja \setin \ObC$;
    \end{enumerate}
    \condit
    \begin{enumerate}
        \item for any object $\Objb \setin \ObC$, there exists a unique morphism of the type $\initobja \to \Objb$.
    \end{enumerate}
\end{ctdefinition}

\begin{example}
    In the category $\CatC = \Set$, the set $\emptyset$ is an initial object.
    Indeed, via the formal definition of functions as special kinds of relations, there is always a unique function from the empty set to any other set.
\end{example}

\begin{example}\label{exa:init-obj-Rel}
    The empty set $\emptyset$ is also an initial object in the category $\Rel$.
    This is an example where a given object is both terminal and initial in the same category.
\end{example}

\todotext{insert more examples}

\begin{gradedexercise}
    Let $\initobja$ and $\initobjb$ be two initial objects in a category $\CatC$.
    Prove that $\initobja$ and $\initobjb$ are isomorphic.
\end{gradedexercise}

